% Figure: intensity transmission coefficient versus impedance ratio
% for a plane wave normally incident on a fluid-fluid interface.
% T_I = 4 r / (1 + r)^2, with r = Z2/Z1. Symmetric about r = 1 in log scale.
\begin{figure}[htbp]
\centering
\begin{tikzpicture}
\begin{axis}[
  width=12cm, height=7.5cm,
  xlabel={impedance ratio $Z_2/Z_1$ (log scale)},
  ylabel={intensity transmission $T_I$},
  xmode=log, log basis x=10,
  xmin=0.001, xmax=1000, ymin=0, ymax=1.1,
  axis lines=left,
  domain=0.001:1000, samples=300,
  every axis x label/.style={at={(current axis.right of origin)}, anchor=west},
  every axis y label/.style={at={(current axis.above origin)}, anchor=south},
  grid=major, grid style={enggray!20},
  tick label style={font=\scriptsize},
  label style={font=\small},
]
% T_I = 4 r / (1 + r)^2
\addplot[engblue, very thick] {4*x/((1+x)^2)};
% perfect match line
\addplot[enggray, thin, dashed] {1};
\node[enggray, font=\scriptsize, anchor=south] at (axis cs:1,1.02) {perfect match $Z_2=Z_1$};
% mark the air-water mismatch at r ~ 3600
\addplot[engred, only marks, mark=*, mark size=2pt] coordinates {(3617,0.0011)};
\node[engred, font=\scriptsize, anchor=west] at (axis cs:30,0.12) {air $\to$ water:};
\node[engred, font=\scriptsize, anchor=west] at (axis cs:30,0.05) {$T_I\approx0.1\%$};
\end{axis}
\end{tikzpicture}
\caption[Intensity transmission versus impedance ratio]{Fraction of
incident intensity transmitted across a normal-incidence fluid-fluid
interface, $T_I = 4r/(1+r)^2$ with $r=Z_2/Z_1$. The curve is symmetric
on the logarithmic ratio axis: a factor-of-ten mismatch transmits the
same fraction whether the wave goes from soft to stiff or stiff to
soft. Transmission peaks at unity only when the impedances match
exactly. The air-to-water case (red marker) sits at $r\approx 3600$,
where only about one part in a thousand of the intensity crosses,
explaining why airborne and waterborne sound are nearly isolated. The
practical lesson is the matching layer: insert a medium of intermediate
impedance to climb the curve in two smaller steps rather than one large
one, exactly the trick a quarter-wave matching layer plays on an
ultrasonic probe.}
\label{fig:vol03ch07:impedance-reflection}
\end{figure}
